\scanpage{32}
plus bas).] \deleted{loc. pour} % L1
\deleted{$\mathcal M^i(X)$ est ``engendré'' par les % L1a
$\mathrm R^i f_{\underline *}$. L'appartenance à $\mathcal M^i(\unsure{X|k(x)})$ \unsure{---} de % L2
$M|_x$ implique l'appart\unsure{en.|enance}} Compatible avec % L3
$f^*$, et \deleted{l'appart} $M|_x\in\mathrm{Ob}\,\mathcal M^{+i}(k(x))$ % L4
$\Longrightarrow\exists\,U\ni x$ voisinage de $\bar x$, tel que % L5
$M|_U\in\mathrm{Ob}\,\mathcal M^{+i}(U)$. Soit $f:X-Y$, % L6
% slanted
\margin{\unsure{généraliser|quasi-limité|quasi-lisse} Compatibilité avec $\varprojlim$ des préschémas ....} % L6a
alors $\mathbb R^j f_!:\mathcal M^{+i}(X)\longrightarrow\mathcal M^{i+j}(Y)$ ; de % L7
plus, $\mathcal M^{+i}(k)$ ($k$ un corps alg. clos) % L8
est ``engendré'' par les \inserted{sous-espaces} $\mathrm R^j f_!(\deleted{\mathbb Z}\mathbb 1_X)$ pour % L9
$X$ si on veut projectif lisse de dim $\leq j$ sur $k$ % L10
[au sens que tout $M\in\mathrm{Ob}\,\mathcal M^i(k)$ a une % L11
filtration dont les facteurs sont isomorphes % L12
à de tels sous-espaces ]. % L13
% slanted
\margin{indices en bas} % L13a
\[
\deleted{\mathcal M^i(X)\otimes\mathbb Z(-1)} % L14
\]
\[
\left\{\begin{aligned}
&\mathcal M^{+i}(X)\otimes\mathcal M^{+j}(X)\subset\mathcal M^{+i+j}(X)\\ % L15
&\mathbb Z(-1)\in\mathrm{Ob}\,\mathcal M^{+\unsure{2|1}}(X) % L16
\end{aligned}\right.
\]
d'où % L17
\[
\mathcal M^{+i}(X)\otimes\mathbb Z(-j)\subset\mathcal M^{+i+2j}(X)\qquad\text{pour }j\geq0 % L18
\]

% archive sheet number: 31
